% chapters/chapter1/1_2_literature.tex
% !TeX root=../../main.tex

\section{پیشینه پژوهش}

کنترل سامانه‌های انتقال بار مانند جرثقیل‌ها به دلیل ماهیت زیرفعال (\lr{Underactuated}) آن‌ها همواره چالشی کلاسیک در مهندسی کنترل بوده است \cite{Ramli2017Review}. مرور جامع پژوهش‌های انجام‌شده در این حوزه نشان می‌دهد که راهبردها عموماً به دو دسته کنترل حلقه-باز (مانند شکل‌دهی ورودی) و کنترل حلقه-بسته (مانند کنترل پیش‌بین، تنظیم‌گر خطی، روش‌های تطبیقی و هوشمند) تفکیک می‌شوند \cite{Ramli2017Review}. با وجود این، ظهور الزامات پیچیده صنعتی (مانند تغییرات متغیرهای فیزیکی در حین کار، دینامیک پیچیده بار و محدودیت‌های عملگرها) موجب شده است تا پژوهشگران به سمت ترکیب روش‌های تحلیلی و الگوریتم‌های بهینه‌سازی سوق پیدا کنند.

در میان راهبردهای حلقه-باز، شکل‌دهی ورودی (\lr{Input Shaping}) کاربرد گسترده‌ای برای سرکوب نوسانات در سامانه‌های انعطاف‌پذیر دارد \cite{Ramli2017Review}. با این وجود، کارایی این روش‌ها در حضور پارامترهای متغیر با زمان، نظیر فرایند بالابر (تغییر طول کابل) و تغییر جرم بار، با افت همراه است \cite{Ramli2018Neural}. برای رفع این محدودیت، رویکردهای تطبیقی توسعه یافته‌اند؛ به عنوان نمونه، تلفیق شبکه‌های عصبی مصنوعی (\lr{ANN}) با الگوریتم‌های بهینه‌سازی ازدحام ذرات (\lr{PSO}) برای تنظیم برخط پارامترهای شکل‌دهنده ورودی با اندازه واحد (\lr{UMZV}) پیشنهاد شده است \cite{Ramli2018Neural}. ارزیابی‌های آزمایشگاهی نشان می‌دهد که این ساختار می‌تواند در حضور عملیات بالابر و تغییرات جرم بار، نوسانات پسماند را تا ۹۱٪ کاهش دهد و در عین حال، زمان پاسخ‌دهی سیستم را نسبت به شکل‌دهنده‌های مقاوم کلاسیک (مانند \lr{ZVDD}) بهبود بخشد \cite{Ramli2018Neural}.

با افزایش پیچیدگی ابعاد فیزیکی و در نظر گرفتن توزیع جرم قلاب و بار، فرض آونگ ساده در بسیاری از کاربردهای صنعتی با خطا همراه است و استفاده از مدل دینامیکی آونگ دوگانه (\lr{Double Pendulum}) ضرورت می‌یابد \cite{Ramli2017Review}. در این دسته از سیستم‌ها، علاوه بر مسئله پایداری، برنامه‌ریزی مسیر با هدف بهینه‌سازی مصرف انرژی و رعایت قیود سخت‌گیرانه حالت (مانند حداکثر سرعت کالسکه و مرزهای مجاز دامنه نوسان) مورد توجه قرار گرفته است \cite{Sun2018Energy}. بر این اساس، راهکارهایی مبتنی بر تبدیل تابع هزینه انرژی به یک مسئله برنامه‌ریزی درجه دوم (\lr{Quadratic Programming}) با استفاده از تکنیک‌های محدب‌سازی (\lr{Convexification}) ارائه شده است \cite{Sun2018Energy}. اجرای این الگوریتم‌ها به صورت برون‌خط (\lr{Off-line}) روی بسترهای آزمایشگاهی آونگ دوگانه ثابت کرد که می‌توان هدف انتقال سریع ارابه و سرکوب کامل نوسانات پسماند را با کمترین مصرف انرژی و بدون نقض قیود سینماتیکی محقق کرد \cite{Sun2018Energy}.

از منظر دینامیک عملگرها و انتقال قدرت، سامانه‌های الکترومکانیکیِ محرکِ ارابه، اغلب تحت تأثیر بارهای اصطکاکی غیرخطی قرار دارند. در مواردی که سیستم در منطقه شیب منفی منحنی اصطکاک کار می‌کند، دینامیک صفرهای سیستم (\lr{Zero Dynamics}) ناپایدار شده و ارتعاشات خودتحریک (\lr{Self-excited Vibrations}) بروز می‌کنند \cite{Golovin2017PFC}. کنترل‌کننده‌های فیدبک کلاسیک در مواجهه با چنین پدیده‌ای دچار محدودیت بهره (\lr{Gain Limitations}) هستند و قادر به پایدارسازی سیستم مرجع نیستند. برای حل این معضل بنیادین، استفاده از جبران‌سازهای پیش‌خورد موازی (\lr{Parallel Feed-forward Compensator}) به عنوان یک افزونه بر ساختار کنترل سنتی پیشنهاد شده است \cite{Golovin2017PFC}. این جبران‌سازها با بازتعریف خروجی و جابه‌جایی صفرهای سیستم به نیمه سمت چپ صفحه مختلط، بستری فراهم می‌کنند تا سیستم الکترومکانیکیِ درگیر با اصطکاک را بتوان با استفاده از روش‌های کنترلی استاندارد و بدون محدودیت بهره، پایدار نمود \cite{Golovin2017PFC}.

مرور این مسیر پژوهشی نشان می‌دهد که اگرچه راهکارهای مجزایی برای بهینه‌سازی مصرف انرژی، انطباق با طول کابل متغیر و رفع ناپایداری عملگرها ارائه شده است، توسعه یک معماری کنترلی جامع که بتواند قیود فیزیکی سخت‌افزار را با دینامیک غیرخطی نوسانات ترکیب کند، همچنان در مرزهای پژوهشی این حوزه قرار دارد.